\chapter{Getting Started}
\label{ch:getting_started_glucat}

\section{Installation}
To install the GluCat C++ library, follow the standard build procedure:

\begin{verbatim}
./configure
make
sudo make install
\end{verbatim}

\subsection{Configuration Options}
You can customize the build using various flags:
\begin{itemize}
    \item \texttt{--with-armadillo}: Use the Armadillo linear algebra library.
    \item \texttt{--with-ordering=boost|std}: Choose the map implementation.
    \item \texttt{--enable-debug=yes|no|full}: Control debugging and optimization levels.
\end{itemize}

\section{Compiling Your First Program}
To compile a C++ program using GluCat, you need to include the header files and link against any necessary libraries (like Armadillo if configured).
\begin{verbatim}
g++ my_program.cpp -o my_program -I/usr/local/include
\end{verbatim}
Refer to \texttt{INSTALL.md} for detailed dependency information.
